\section{Background and Preliminaries}
\label{sec:background}

Speculative decoding accelerates autoregressive generation by pairing a small
draft model with a large target verifier \citep{leviathan2023speculative}. Let
$q_i(\cdot \mid h_i)$ denote the draft distribution at speculative step $i$
given history $h_i$, and let $p_i(\cdot \mid h_i)$ denote the corresponding
target distribution. The draft proposes a token $y_i \sim q_i$, and the target
accepts it with probability
\[
\alpha_i = \min\!\left(1, \frac{p_i(y_i)}{q_i(y_i)}\right).
\]

The verifier either commits the drafted token or replaces it with a
target-sampled correction. Consequently, speculative efficiency depends
directly on the overlap between $p_i$ and $q_i$: high agreement yields long
accepted draft blocks, whereas low agreement forces frequent verifier
intervention.

This dependence can be written in a particularly convenient form. Under draft
sampling, the one-step expected acceptance satisfies
\[
\mathbb{E}[\alpha_i] = 1 - \tv(p_i, q_i),
\]
where $\tv$ is total variation distance. This identity makes the attack surface
explicit. Any prompt that increases early draft-target disagreement will, in
expectation, reduce verifier acceptance and therefore diminish speculative
speedup.

The systems metrics used throughout the paper are derived from the same
acceptance process. Let accepted drafted tokens denote the number of draft
tokens that survive verification, drafted tokens denote the total number of
tokens proposed by the draft, and target steps denote the number of verifier
steps taken by the target. Then
\[
\text{aggregate acceptance}
= \frac{\text{accepted drafted tokens}}{\text{drafted tokens}},
\]
\[
\text{block efficiency}
= \frac{\text{accepted drafted tokens}}{\text{target steps}}.
\]
Aggregate acceptance captures how much of the draft computation survives, while
block efficiency measures how much useful drafted work is recovered per
verifier step. These counter-based metrics complement wall-clock latency in the
experiments, and they are less sensitive to incidental runtime noise than raw
timing alone.
